\documentclass[12pt]{article}
\usepackage[utf8]{inputenc}

%%% I added a few packages. This next one allows you to set the margins.
\usepackage[top=.5in,bottom=1in,right=.75in,left=.75in]{geometry}
\usepackage{multicol}
%\usepackage{newspaper}
\setlength{\columnsep}{1cm}
\usepackage{blindtext}
\usepackage{varwidth}

%%% This one does boxes in a really nice way.
\usepackage{tcolorbox}
\usepackage{mdframed}
\usepackage{minibox}
\usepackage{graphicx}
\graphicspath{ {./images/} }
\usepackage{epstopdf}
  \epstopdfDeclareGraphicsRule{.pdf}{png}{.png}{convert #1 \OutputFile}
  \DeclareGraphicsExtensions{.png,.pdf}
\usepackage{hyperref}
\hypersetup{
    colorlinks=true,
    linkcolor=blue,
    filecolor=magenta,      
    urlcolor=cyan,
    pdftitle={MSCSMess},
    pdfpagemode=FullScreen,
    }
\urlstyle{same}
 \usepackage{fix-cm}
 \usepackage{soul}
 \usepackage{graphicx}
 \usepackage{multicol}


\newcommand{\ds}{\displaystyle}
\newcommand{\ZZ}{\mathbb{Z}}
\newcommand{\QQ}{\mathbb{Q}}
\newcommand{\RR}{\mathbb{R}}
\newcommand{\CC}{\mathbb{C}}
\newcommand{\FF}{\mathbb{F}}
\newcommand{\xx}{\mathbf{x}}
\newcommand{\uu}{\mathbf{u}}
\newcommand{\vv}{\mathbf{v}}

\newcommand{\bolder}{\textbf}
 
 
\title{\fontsize{60}{70}\selectfont {\bf {MSCS} 
\includegraphics[scale=.3]{images/OleLion.png} {Mess}}}
\author{Department of Mathematics, Statistics, and Computer Science \\ St.\@ Olaf College, Northfield, MN 55057 \vspace{-2ex}} 
\date{\today\, $|$ Volume 54, No.\@ 1  \line(1,0){500}}


\begin{document}

\maketitle


\begin{center}
{\bf Welcome to Volume 54 of the MSCS Mess!}
\end{center}
We hope that your first week of classes has gone as smoothly as you hoped. If you haven't read the Mess before, the MSCS Mess is our weekly departmental newsletter, and it's all about \bolder{informing you about relevant events for Math, Statistics, Data Science, and Computer Science students}. These events include Colloquium talks, Research Seminars, celebratory events, and some fun facts/applications. To submit anything for publication, email \textit{messeditor@stolaf.edu}.


 \begin{multicols}{2}

\begin{tcolorbox}[colback=white, colframe=black]
    \begin{center}
    \title{{\large \bolder{MSCS Showcase}}}
    \end{center}
\end{tcolorbox}

Are you a newcomer to MSCS looking to see what we have to offer? Are you an experienced member of MSCS looking to get more involved? Come to the Taylor Swift-inspired  MSCS Showcase on Thursday, September 25 at 4:30 pm in the Buntrock Commons Ballrooms! The MSCS faculty will talk about the many things that make our department special – like classes, research opportunities, and more!
 

\begin{tcolorbox}[colback=white, colframe=black]
    \begin{center}
    \title{{\large \bolder{Applications open for Math 257, J-term 2026}}}
    \end{center}
\end{tcolorbox}

Math 257 is a course on the mathematics of Emmy Noether and Sofia Kovalevskaya and a discussion of equitable access to education which spends 3 weeks in Germany and 1 week in Sweden. Applications for this coming J-term are being accepted through \textbf{September 18}. The course counts toward the math major, its only prerequisite is a proof-writing course (current registration in Math 244 or 252 counts), and it is super fun! Ask Prof. Dietz for details.





\begin{tcolorbox}[colback=white, colframe=black]
    \begin{center}
    \title{{\large \bolder{Problem Solving Group}}}
    \end{center}
\end{tcolorbox}
The Math Problem Solving Group meets weekly to have fun, solve math problems, eat candy, and prepare for upcoming math contests. Meetings are 11:30am-12:30pm, in RNS 207 (the Center for Interdisciplinary Research room). Everyone is welcome, regardless of prior experience in mathematics or math contests.

The first contest will be the Jerry Heuer Team Contest on the morning of Saturday, November 15. It's a fun, accessible contest for teams of up to three students. For more information, or to be added to the problem solving group email list, contact Prof. Wright at wright5@stolaf.edu.


\begin{tcolorbox}[colback=white, colframe=black]
    \begin{center}
    \title{{\large \bolder{CURI Fall Showcase}}}
    \end{center}
\end{tcolorbox}

On Friday, October 10, 4-6pm, in the Kings Dining Room Complex, the Collaborative Undergraduate Research \& Inquiry (CURI) program will sponsor a poster session that includes all of the summer CURI researchers, Steen Fellows, Magnus the Good Fellows, and academic year researchers. There will be several projects from MSCS and a lot from the sciences in general. Support your friends and peers, and eat some snacks!




\end{multicols}

\vfill

\begin{center}

\includegraphics[scale=0.45]{images/2025 MSCS Showcase.png}
\end{center}

\vspace{0.3in}
\noindent{\it 

\begin{center}


\includegraphics[scale=0.8]{images/T-Shirt Design Contest Flyer.png}
\end{center}

\vspace{0.3in}
\noindent{\it 


\begin{center}

\includegraphics[scale=0.55]{images/Colloquium Kelsey Grinde.png}
\end{center}

\vspace{0.3in}
\noindent{\it 

\begin{center}

\includegraphics[scale=0.8]{images/AWM CEO Seminar Series - National Labs.png}
\end{center}

\vspace{0.3in}
\noindent{\it 

\begin{center}

\includegraphics[scale=0.7]{images/Tea Time Poster.png}
\end{center}

\vspace{0.3in}
\noindent{\it 



\noindent{To submit an article, event, or anything else for publication in the Mess, email messeditor@stolaf.edu. Visit \href{https://wp.stolaf.edu/mscs/mscs-mess/}{http://wp.stolaf.edu/mscs/mscs-mess/} for a digital archive of previous MSCS Mess issues.}}\\
	
\hfill Nima Dahir, Editor
	
\hfill Kyle Salois, Faculty Advisor


 

\end{document}
